\section*{Summary}
\addcontentsline{toc}{section}{Summary}

This document develops an improved electoral procedure for the Bundestag based on three building blocks:

\begin{enumerate}
  \item \textbf{Iterative Proportional Representation} --- Successive rounds of voting ensure
    that there is always an election winner who has been elected by at least half of the people.
    Proportionality is limited to the opposition. Coalitions remain possible,
    but governmental capability is guaranteed even without them.

  \item \textbf{Direct Election of Lead Candidates} --- Each party runs with three
    lead candidates. Every voter casts a \enquote{party vote} (which determines the
    seat distribution in parliament) and a \enquote{lead candidate vote}. With the
    lead candidate vote, voters select their party's lead candidate.

  \item \textbf{Personal Representation of Each Constituency} --- Voters can no longer directly
    determine which party should represent their constituency. This is determined by an algorithm
    that calculates the relative importance of the constituency for each party based on the party vote
    distribution and assigns constituencies to parties in a way that maximizes the total relative
    importance.
    Each party must nominate three candidate representatives at the constituency level. Voters
    can select the most suitable representative for every party in their constituency (not just their own).
    The candidate representative of the assigned party with the most votes wins the constituency.
\end{enumerate}

An open-source simulation of the procedure is available as a
\href{https://github.com/danielschwalbe/ipres}{GitHub project}.
